\documentclass[portuguese,12pt,a4paper,onecolumn]{article}
\usepackage[T1]{fontenc}
\usepackage[left=3cm, right=3cm, top=2cm, bottom=2cm]{geometry}
\usepackage{graphicx}
\usepackage{amssymb}
\usepackage{amsthm}
\usepackage{nameref}
\usepackage[portuguese]{babel}
\usepackage{hyperref}
\hypersetup{
	colorlinks  = true, %Colours links instead of ugly boxes
	linkcolor = blue,
	urlcolor = blue,
	citecolor = black,
}
\usepackage{siunitx}
\usepackage{booktabs}
\usepackage{graphicx,xcolor}
\sisetup{	locale = FR,
	text-series-to-math = true ,
	propagate-math-font = true
}
\usepackage[version=4]{mhchem}
\usepackage{chemmacros}
\usepackage{chemfig}
\chemsetup{
	formula=mhchem,
	reactions/tag-open=(R,
	reactions/tag-close=),
}
\renewcommand\reactionautorefname{Reação}
\usepackage{listings}
% Configure style
\lstset{
	backgroundcolor=\color{gray!10},
	basicstyle=\ttfamily\footnotesize,
	breaklines=true,
	numbers=left,
	numberstyle=\tiny,
	commentstyle=\color{green},
	keywordstyle=\color{blue},
	stringstyle=\color{purple}
}
\renewcommand{\lstlistingname}{Código-fonte}

\title{Conversão do modelo distribuído de Síntese de Fischer-Tropsch para parâmetros concentrados}
\author{Rafael Belo Duarte}
\begin{document}
	\maketitle
	\section{Resumo do modelo de parâmetros distribuídos}
	
	Antes de mostrar o modelo de parâmetros concentrados (\autoref{lumped_model}) vamos resumir o de parâmetros distribuídos.
	
	\subsection{Reação de deslocamento}
	
	\begin{reaction}
		CO + H2O <=> CO2 + H2
		\label{reac:WGS}
	\end{reaction}
	
	\begin{equation}
		r_\mathrm{WGS}' = k_{\mathrm{WGS}}' P_{\ce{CO}}^a P_{\ce{H2O}}^b  \left(1 - \frac{1}{K_{\mathrm{WGS}}} \frac{P_{\ce{CO2}} P_{\ce{H2}}}{P_{\ce{CO}} P_{\ce{H2O}}} \right)
		\label{eq:rate_CO2}
	\end{equation}
	
	\begin{equation}
		k'_\mathrm{WGS} = A_\mathrm{WGS}' \exp \left( \frac{-E_\mathrm{WGS}}{RT} \right)
		\label{eq:k_WGS}
	\end{equation}
	
	\begin{equation}
		K_{\mathrm{WGS}} = \num{1.45e-2} \exp \left( \frac{\num{4.62e3}}{T} \right)
		\label{eq:K_WGS}
	\end{equation}
	
	\begin{table*}[!h]
		\centering
		\caption{Parâmetros da reação de deslocamento.}
		\label{tab:params_WGS}
		%	\resizebox{1\textwidth}{!}{%
			\begin{tabular}{rrrr}
				\toprule
				$A_{\mathrm{WGS}}'$ (\unit{\mole \per \second \per \gram \per \pascal^{a+b}}) & $E_{\mathrm{WGS}}$ (\unit{\joule \per \mole}) & $a$ & $b$ \\
				\midrule
				\num{2.066E+06} & \num{164536} & \num{0.02} & \num{0.53} \\
				\bottomrule
			\end{tabular}%}
	\end{table*}
	
	\subsection{Reações de síntese}
	
	\subsubsection{Reações}
	
	\paragraph{Parafinas}
	\begin{reaction}
		n CO + $(2n + 1)$ H2 -> C_n H_{2n+2} + n H2O
	\end{reaction}
	
	\paragraph{Olefinas}
	\begin{reaction}
		n CO + $2 n$ H2 -> C_n H_{2n} + n H2O
	\end{reaction}
	
	\subsubsection{Velocidades de reação}
	
	Nas equações das velocidades de reação $n$ é o comprimento de cadeia, $y_n$ (\autoref{eq:alpha}) a fração molar de hidrocarbonetos de tamanho de cadeia $n$ (excluídos os demais componentes: \ce{H2O}, \ce{CO2}, \ce{CO} e \ce{H2}), $r_\mathrm{HCs}'$ é velocidade global de formação de hidrocarbonetos (\autoref{eq:r_HCs}), e $\upsilon_n$ a probabilidade de hidrogenação (\autoref{eq:upsilon_n_logistic}).
	
	\paragraph{Parafinas}
	\begin{equation}
		r_{\mathrm{parafina},n} = y_n \upsilon_n r_{\mathrm{HCs}}'
	\end{equation}
	
	\paragraph{Olefinas}
	\begin{equation}
		r_{\mathrm{olefina},n} = y_n (1 - \upsilon_n) r_{\mathrm{HCs}}'
	\end{equation}
	
	\paragraph{Monóxido de carbono}
	\begin{equation}
		- r_{\ce{CO}}' = r_{\mathrm{WGS}}' + \sum_{n=1}^{n \geq 30} r_{\mathrm{HCs}}' n y_n
	\end{equation}
	
	\paragraph{Hidrogênio}
	\begin{equation}
		-r_{\ce{H2}}' = \left(\sum_{n = 1}^{N} r_{\mathrm{HCs}}' [y_n (2n + \upsilon_n)]\right) - r_{\mathrm{WGS}}'
		\label{eq:r_H2_simple}
	\end{equation}
	
	\paragraph{Água}
	\begin{equation}
		r_{\ce{H2O}}' = \left(\sum_{n=1}^{n \geq 30} r_{\mathrm{HCs}}' n y_n\right) - r_{\mathrm{WGS}}'
	\end{equation}
	
	\paragraph{Dióxido de carbono}
	\begin{equation}
		r_{\ce{CO2}}' = r_{\mathrm{WGS}}'
	\end{equation}
	
	\paragraph{A velocidade de formação dos hidrocarbonetos}
	
	\begin{equation}
		r_\mathrm{HCs}' = k' \frac{P_{\ce{H2}}^{0.75} P_{\ce{CO}}}{(1 + K_{\ce{CO}} P_{\ce{CO}})^2}
		\label{eq:r_HCs}
	\end{equation}
	
	\begin{equation}
		k' = A_\mathrm{HCs}' \exp \left( \frac{-E_\mathrm{HCs}}{RT} \right)
		\label{eq:k}
	\end{equation}
	
	\begin{equation}
		K_{\ce{CO}} = k_{\ce{CO}} \exp \left( \frac{-\Delta H_{\ce{CO}}}{RT} \right)
		\label{eq:K_CO}
	\end{equation}
	
	\colorbox{pink}{\parbox{\textwidth}{
			\textbf{\underline{IMPORTANTE}}: Para os hidrocarbonetos com $n = 1$ e $n = 2$ substitua a energia de ativação da \autoref{eq:k} ($E_\mathrm{HCs}$) por $E_{\mathrm{C1}}$ e $E_{\mathrm{C2}}$ (\autoref{tab:params_Mousavi}), respectivamente.
		}}
	
	\begin{table*}[!h]
		\centering
		\caption{Parâmetros das Equações \ref{eq:r_HCs}, \ref{eq:k} e \ref{eq:K_CO}.}
		\label{tab:params_Mousavi}
		\resizebox{1\textwidth}{!}{%
			\begin{tabular}{rrrrrr}
				\toprule
				$k_{\ce{CO}}$ (\unit{\pascal^{-1}}) & $\Delta H_{\ce{CO}}$ (\unit{\joule \per \mole}) & $A_{\mathrm{HCs}}'$ (\unit{\mole \per \second \per \gram \per \pascal^{1.75}}) & $E_{\mathrm{HCs}}$ (\unit{\joule \per \mole}) & $E_{\mathrm{C1}}$ (\unit{\joule \per \mole}) & $E_{\mathrm{C2}}$ (\unit{\joule \per \mole}) \\
				\midrule
				\num{1.398E-05} & \num{933.0} & \num{9.551E-01} & \num{147204} & \num{142602} & \num{153273} \\
				\bottomrule
			\end{tabular}
		}
	\end{table*}
		
	\subsubsection{Distribuição dos hidrocarbonetos}
	
	\begin{equation}
		y_n = (1 - \alpha) \alpha^{(n - 1)}
		\label{eq:alpha}
	\end{equation}
	
	$y_n$ é a fração molar do hidrocarboneto com $n$ carbonos. O $y_n$ representa exclusivamente a distribuição dos hidrocarbonetos, não estão inclusos o \ce{CO}, \ce{CO2}, \ce{H2} e \ce{H2O}.
	
	A probabilidade de crescimento de cadeia ($\alpha$) é:
	
	\begin{equation}
		\alpha = \hat{\beta_0} + \hat{\beta_1} T
	\end{equation}

	Onde $T$ está em Kelvin e [$\hat{\beta_0}, \hat{\beta_1}] = [\num{1.626}, \qty{-1.495e-03}{\per \kelvin}]$.
	
	\subsubsection{Probabilidade de hidrogenação}
	
	\begin{equation}
		\upsilon_n = \frac{1}{1 + [1 - (n \odot 1) - (n \odot 3)] \mathrm{exp}[-(\hat{\beta}_3 + \hat{\beta}_4 n + \hat{\beta}_5 T + \hat{\beta}_6 (n \odot 2))]}
		\label{eq:upsilon_n_logistic}
	\end{equation}
	
	Onde [$\hat{\beta}_3, \hat{\beta}_4, \hat{\beta}_5, \hat{\beta}_6$] = [-13,56; 0,3305; \qty{0,02179}{\per \kelvin}; 2,291]. O símbolo $\odot$ representa o operador XNOR, que retorna True (1) somente quando ambas as entradas são idênticas. O termo $[1 - (n \odot 1) - (n \odot 3)]$ serve para forçar $\upsilon = 1$ para metano e propano. Essa expressão permite que o modelo descreva a proporção de etano e propano maior do que a esperada em comparação com o restante da distribuição. A \autoref{eq:upsilon_n_logistic} pode ser inserida em um simulador por meio da criação de uma planilha ou usando \textit{scripts} internos, uma funcionalidade presente tanto no DWSIM quanto no Aspen Hysys.
	
	\textbf{No modelo de parâmetros concentrados esta equação não será necessária}, pois estaremos interessados apenas em estimar a produção de hidrocarbonetos por tamanho de cadeia, descartando esta distinção entre hidrocarbonetos saturados e insaturados.
	
	\newpage
	
	\section{O modelo de parâmetros concentrados}
	\label{lumped_model}
	
	\subsection{Definição das frações do produto}
	
	Para criar o modelo de parâmetros concentrados, os produtos foram agrupados por tamanho de cadeia de acordo com a \autoref{tab:fracoes}.
	
	\begin{table*}[!h]
		\centering
		\caption{Frações do produto no modelo de parâmetros concentrados.}
		\label{tab:fracoes}
		\begin{tabular}{cc}
			\hline
			Produto & Tamanho de cadeia ($n$) \\
			\hline
			Gases & 1 - 4 \\
			Gasolina & 5 - 10 \\
			Querosene & 11 - 16 \\
			Diesel & 17 - 20 \\
			Parafinas pesadas & 21 + \\
			\hline
		\end{tabular}
	\end{table*}
	
	\subsection{Criação de pseudo-componentes no DWSIM}
	
	Para agrupar os componentes vamos primeiro realizar uma simulação usando o modelo de parâmetros distribuídos, separando os produtos por tamanho de cadeia. Destas correntes criamos os pseudo-componentes que serão usados para validar o modelo de parâmetros concentrados.
	
	Três informações são necessárias para criar os pseudo-componentes: a massa molar, a densidade relativa a \qty{60}{\degree F} (\qty{15,5556}{\degreeCelsius}) e o ponto de bolha, contidos na \autoref{tab:simul_pseudo_comps}.
	
	\begin{table*}[!h]
		\centering
		\caption{Resultados de uma simulação usando o modelo de parâmetros distribuídos a \qty{20}{\bar}, \qty{240}{\degreeCelsius} e com \ce{H2}/\ce{CO} = 2.}
		\label{tab:simul_pseudo_comps}
		\resizebox{\textwidth}{!}{%
			\begin{tabular}{lllllll}
				\hline
				Pseudo-componente         & querosene   & parafinas pesadas & gasolina   & gases       & diesel      & Unidade \\ \hline
				Temperatura               & 15,5556     & 15,5556           & 15,5556    & 15,5556     & 15,5556     & \unit{\degreeCelsius}     \\
				Pressão                   & 1           & 1                 & 1          & 1           & 1           & atm   \\
				Vazão mássica             & 0,00422758  & 0,00239809        & 0,00510504 & 0,00267842  & 0,00160822  & \unit{\kilogram~d^{-1}}  \\
				Vazão molar               & 2,69644E-07 & 7,79018E-08       & 5,9037E-07 & 1,15091E-06 & 7,18772E-08 & \unit{\mol \per \second} \\
				Vazão Volumétrica         & 3,87331E-06 & 2,20618E-06       & 5,1551E-06 & 0,00163585  & 1,42698E-06 & \unit{\liter \per \minute} \\
				Densidade da mistura      & 0,757963    & 0,754853          & 0,687701   & 0,00113703  & 0,782641    & \unit{\gram \per \centi \meter \cubed} \\
				Peso molecular da mistura & 181,463     & 356,29            & 100,083    & 26,9353     & 258,964     & \unit{\gram \per \mol} \\
				\begin{tabular}[c]{@{}l@{}}Ponto de bolha \\ na pressão da corrente\end{tabular} & 220,807 & 385,251 & 67,024 & -156,049 & 318,332 & \unit{\degreeCelsius} \\ \hline
			\end{tabular}%
		}
	\end{table*}
	
	Para os componentes líquidos foi possível criar pseudo-componentes sem problema, mas para os gases o algoritmo retorna um pseudo-componente com fórmula \ce{C0H0}. Portanto, não vamos usar este pseudo-componente nas simulações, em seu lugar agruparemos as velocidades de reação do \ce{C2} ao \ce{C4} dentro da reação de formação do metano. As propriedades dos demais pseudo-componentes são fornecidas  na \autoref{tab:pseudo_props}.
	
	% Please add the following required packages to your document preamble:
	% \usepackage{graphicx}
	\begin{table}[!h]
		\centering
		\caption{Propriedades dos pseudo-componentes.}
		\label{tab:pseudo_props}
		\resizebox{\textwidth}{!}{%
			\begin{tabular}{llllll}
				\hline
				BASIC PROPERTIES &
				&
				&
				&
				&
				\\ \hline
				Name &
				gasolina &
				querosene &
				diesel &
				pesados &
				Unidades \\ \hline
				Database &
				DWSIM &
				DWSIM &
				DWSIM &
				DWSIM &
				\\
				Type &
				\begin{tabular}[c]{@{}l@{}}Petroleum Fraction \\ (Pseudocompound)\end{tabular} &
				\begin{tabular}[c]{@{}l@{}}Petroleum Fraction \\ (Pseudocompound)\end{tabular} &
				\begin{tabular}[c]{@{}l@{}}Petroleum Fraction \\ (Pseudocompound)\end{tabular} &
				\begin{tabular}[c]{@{}l@{}}Petroleum Fraction \\ (Pseudocompound)\end{tabular} &
				\\
				CAS ID &
				&
				&
				&
				&
				\\
				Molecular Weight &
				100,083 &
				181,463 &
				258,964 &
				356,29 &
				\\
				Critical Temperature &
				512,247151116108 &
				667,094185221083 &
				749,407462353205 &
				781,834406669311 &
				K \\
				Critical Pressure &
				3218488,97808149 &
				1729493,86407653 &
				1164581,37918558 &
				750488,251805917 &
				Pa \\
				Critical Volume &
				0,35700246352383 &
				0,79940946941013 &
				1,23826310774281 &
				1,77741661629187 &
				\unit{\meter \cubed \per \kilo \mol} \\
				Critical Compressibility &
				0,269794857423962 &
				0,249282125304351 &
				0,231448698092234 &
				0,205214637636691 &
				\\
				Acentric Factor &
				0,265064282200474 &
				0,521473433695612 &
				0,744391273847072 &
				1,07231702954136 &
				\\
				\begin{tabular}[c]{@{}l@{}}Gibbs Energy of Formation \\ (Ideal Gas at 298.15 K)\end{tabular} &
				-1179,11352640155 &
				-979,104077152584 &
				-923,137997948034 &
				-916,910482319602 &
				\unit{\kilo \joule \per \kilogram} \\
				\begin{tabular}[c]{@{}l@{}}Enthalpy of Formation \\ (Ideal Gas at 298.15 K)\end{tabular} &
				-2,1028883108762 &
				-1,8421528395773 &
				-1,85610400388881 &
				-1,26097465419024 &
				\unit{\kilo \joule \per \kilogram} \\
				Normal Boiling Point &
				340,174 &
				493,957 &
				591,482 &
				658,401 &
				K \\
				&
				&
				&
				&
				&
				\\ \hline
				\begin{tabular}[c]{@{}l@{}}MODEL-SPECIFIC \\ PROPERTIES\end{tabular} &
				&
				&
				&
				&
				\\ \hline
				\begin{tabular}[c]{@{}l@{}}Chao Seader \\ Acentric Factor\end{tabular} &
				0,265064282200474 &
				0,521473433695612 &
				0,744391273847072 &
				1,07231702954136 &
				\\
				\begin{tabular}[c]{@{}l@{}}Chao Seader \\ Solubility Parameter\end{tabular} &
				6,74732726610346 &
				5,51637798217385 &
				4,8923126589583 &
				4,3130662093866 &
				\\
				\begin{tabular}[c]{@{}l@{}}Chao Seader \\ Liquid Molar Volume\end{tabular} &
				134,61042873731 &
				301,417198396027 &
				463,576964325648 &
				651,342364963486 &
				\\
				Rackett Compressibility &
				0,269794857423962 &
				0,249282125304351 &
				0,231448698092234 &
				0,205214637636691 &
				\\
				&
				&
				&
				&
				&
				\\ \hline
				\begin{tabular}[c]{@{}l@{}}PSEUDOCOMPOUND \\ (PETROLEUM FRACTION)\\ RELATED PROPERTIES\end{tabular} &
				&
				&
				&
				&
				\\ \hline
				Specific Gravity &
				0,687701 &
				0,757963 &
				0,782641 &
				0,754853 &
				\\
				Watson K &
				12,3478558667387 &
				12,686423659089 &
				13,0469425638431 &
				14,0192660269206 &
				\\
				\begin{tabular}[c]{@{}l@{}}Temperature for \\ Viscosity Data Point 1\end{tabular} &
				310,95 &
				310,95 &
				310,95 &
				310,95 &
				K \\
				Viscosity @ T1 &
				2,90960390372366E-07 &
				9,58036292121748E-07 &
				2,72557478768417E-06 &
				4,73752008031969E-06 &
				\unit{\pascal \second} \\
				\begin{tabular}[c]{@{}l@{}}Temperature for \\ Viscosity Data Point 2\end{tabular} &
				372,05 &
				372,05 &
				372,05 &
				372,05 &
				K \\
				Viscosity @ T2 &
				3,01768153261329E-07 &
				6,15640578277514E-07 &
				1,29660618472763E-06 &
				1,60355612823309E-06 &
				\unit{\pascal \second} \\ \hline
			\end{tabular}%
		}
	\end{table}
	
	\subsection{Reação de deslocamento}
	
	A reação de deslocamento não muda no modelo de parâmetros concentrados.
	
	\begin{reaction}
		CO + H2O <=> CO2 + H2
		\label{reac:WGS_2}
	\end{reaction}
	
	\begin{equation}
		r_\mathrm{WGS}' = k_{\mathrm{WGS}}' P_{\ce{CO}}^a P_{\ce{H2O}}^b  \left(1 - \frac{1}{K_{\mathrm{WGS}}} \frac{P_{\ce{CO2}} P_{\ce{H2}}}{P_{\ce{CO}} P_{\ce{H2O}}} \right)
		\label{eq:rate_CO2_2}
	\end{equation}
	
	\begin{equation}
		k'_\mathrm{WGS} = A_\mathrm{WGS}' \exp \left( \frac{-E_\mathrm{WGS}}{RT} \right)
		\label{eq:k_WGS_2}
	\end{equation}
	
	\begin{equation}
		K_{\mathrm{WGS}} = \num{1.45e-2} \exp \left( \frac{\num{4.62e3}}{T} \right)
		\label{eq:K_WGS_2}
	\end{equation}
	
	\begin{table*}[!h]
		\centering
		\caption{Parâmetros da reação de deslocamento.}
		\label{tab:params_WGS_2}
		%	\resizebox{1\textwidth}{!}{%
			\begin{tabular}{rrrr}
				\toprule
				$A_{\mathrm{WGS}}'$ (\unit{\mole \per \second \per \gram \per \pascal^{a+b}}) & $E_{\mathrm{WGS}}$ (\unit{\joule \per \mole}) & $a$ & $b$ \\
				\midrule
				\num{2.066E+06} & \num{164536} & \num{0.02} & \num{0.53} \\
				\bottomrule
			\end{tabular}%}
	\end{table*}
	
	\subsection{Reações de formação dos pseudo-componentes}
	
	As fórmulas moleculares dos nossos pseudo-componentes são:

	\begin{itemize}
		\item Gases: \ce{CH4}
		\item Gasolina: \ce{C_{6,95}H_{15,50}}
		\item Querosene: \ce{C_{12,78}H_{27,32}}
		\item Diesel: \ce{C_{18,26}H_{37,09}}
		\item Pesados: \ce{C_{24,94}H_{48,14}}
	\end{itemize}
	
	E suas reações de formação:
	
	\paragraph{Gases}
	\begin{reaction}
		3 H2 + CO -> CH4 + H2O
		\label{reac:gases_2}
	\end{reaction}
	
	\paragraph{Gasolina}
	\begin{reaction}
		14,9 H2 + 6,95 CO -> 6,91768 H2O + C_{6,95}H_{15,50}
	\end{reaction}
	
	\paragraph{Querosene}
	\begin{reaction}
		26,56 H2 + 12,78 CO -> 12,7696 H2O + C_{12,78}H_{27,32}
		\label{reac:querosene}
	\end{reaction}
	
	\paragraph{Diesel}
	\begin{reaction}
		37,52 H2 + 18,26 CO -> 18,2143 H2O + C_{18,26}H_{37,09} 
	\end{reaction}
	
	\paragraph{Pesados}
	\begin{reaction}
		50,88 H2 + 24,94 CO -> 24,6929 H2O + C_{24,94}H_{48,14}
		\label{reac:pesados_2}
	\end{reaction}
	
	\colorbox{yellow}{\parbox{\textwidth}{
			\textbf{\underline{CUIDADO}}: As relações estequiométricas acima são apenas aproximações, os balanços molares não vão fechar exatamente pois os pseudo-componentes originam de duas famílias de reações: de formação de parafinas e de olefinas.
	}}
	
	\subsection{Velocidades de reação do modelo de parâmetros concentrados}
	
	Para obter as velocidades de reação dos pseudo-componentes, podemos somar as velocidades de reação dos componentes individuais do modelo distribuído:
	
	\paragraph{Gases}
	\begin{equation}
		r'_\mathrm{gases} = \sum_{n=1}^{n=4} r'_\mathrm{ HCs } y_n
	\end{equation}
	
	\colorbox{pink}{\parbox{\textwidth}{
			\textbf{\underline{IMPORTANTE}}: Para os hidrocarbonetos com $n = 1$ e $n = 2$ substitua a energia de ativação da \autoref{eq:k_2} ($E_\mathrm{HCs}$) por $E_{\mathrm{C1}}$ e $E_{\mathrm{C2}}$ (\autoref{tab:params_Mousavi_2}), respectivamente.
	}}

	Ou seja:
	\begin{equation}
		r'_\mathrm{gases} = r_1 + r_2 + r_3 + r_4
		\label{eq:r_gases_2}
	\end{equation}
	
	Com:
	\begin{equation}
		r_1 = A'_\mathrm{ HCs } \exp \left ( \frac{-E_\mathrm{C1}}{RT} \right ) \frac{P_{\ce{H2}}^{0.75}P_\mathrm{CO}}{(1 + K_\mathrm{CO}P_\mathrm{CO})^2} y_1
	\end{equation}
	
	\begin{equation}
		r_2 = A'_\mathrm{ HCs } \exp \left ( \frac{-E_\mathrm{C2}}{RT} \right ) \frac{P_{\ce{H2}}^{0.75}P_\mathrm{CO}}{(1 + K_\mathrm{CO}P_\mathrm{CO})^2} y_2
	\end{equation}
	
	\begin{equation}
		r_3 = A'_\mathrm{ HCs } \exp \left ( \frac{-E_\mathrm{HCs}}{RT} \right ) \frac{P_{\ce{H2}}^{0.75}P_\mathrm{CO}}{(1 + K_\mathrm{CO}P_\mathrm{CO})^2} y_3
	\end{equation}
	
	\begin{equation}
		r_4 = A'_\mathrm{ HCs } \exp \left ( \frac{-E_\mathrm{HCs}}{RT} \right ) \frac{P_{\ce{H2}}^{0.75}P_\mathrm{CO}}{(1 + K_\mathrm{CO}P_\mathrm{CO})^2} y_4
	\end{equation}
	

	\paragraph{Gasolina}
	\begin{equation}
		r'_\mathrm{gasolina} = r'_\mathrm{ HCs } \sum_{n=5}^{n=10} y_n
	\end{equation}
	
	\paragraph{Querosene}
	\begin{equation}
		r'_\mathrm{querosene} = r'_\mathrm{ HCs } \sum_{n=11}^{n=16} y_n
	\end{equation}
	
	\paragraph{Diesel}
	\begin{equation}
		r'_\mathrm{diesel} = r'_\mathrm{ HCs } \sum_{n=17}^{n=20} y_n
	\end{equation}
	
	\paragraph{Pesados}
	\begin{equation}
		r'_\mathrm{pesados} = r'_\mathrm{ HCs } \sum_{n=21}^{n \geq 30} y_n
	\end{equation}
	
	Veja que nas equações acima o $n$ não multiplica o $y_n$ pois estamos seguindo as estequiometrias das Reações \ref{reac:gases_2} a \ref{reac:pesados_2}. Enquanto para o \ce{CO}, na \autoref{eq:r_monox}, o $n$ aparece pois a cada mol de hidrocarboneto que é gerado $n$ mols de monóxido de carbono são consumidos, a lógica é a mesma para o consumo de hidrogênio e a geração de água.
	
	\paragraph{Monóxido de carbono}
	\begin{equation}
		- r_{\ce{CO}}' = r_{\mathrm{WGS}}' + \sum_{n=1}^{n \geq 30} r_{\mathrm{HCs}}' n y_n
		\label{eq:r_monox}
	\end{equation}
	
	\paragraph{Hidrogênio}
	\begin{equation}
		-r_{\ce{H2}}' = 3 r'_\mathrm{gases} + 14,9 r'_\mathrm{gasolina} + 26,56 r'_\mathrm{querosene} + 37,52 r'_\mathrm{diesel} + 50,88 r'_\mathrm{pesados}
		\label{eq:r_H2_simple_2}
	\end{equation}
	
	\paragraph{Água}
	\begin{equation}
		r_{\ce{H2O}}' =  \left(\sum_{n=1}^{n \geq 30} r_{\mathrm{HCs}}' n y_n\right) - r_{\mathrm{WGS}}'
	\end{equation}
	
	\paragraph{Dióxido de carbono}
	\begin{equation}
		r_{\ce{CO2}}' = r_{\mathrm{WGS}}'
	\end{equation}
	
	\paragraph{A velocidade de formação dos hidrocarbonetos}
	
	\begin{equation}
		r_\mathrm{HCs}' = k' \frac{P_{\ce{H2}}^{0.75} P_{\ce{CO}}}{(1 + K_{\ce{CO}} P_{\ce{CO}})^2}
		\label{eq:r_HCs_2}
	\end{equation}
	
	\begin{equation}
		k' = A_\mathrm{HCs}' \exp \left( \frac{-E_\mathrm{HCs}}{RT} \right)
		\label{eq:k_2}
	\end{equation}
	
	\begin{equation}
		K_{\ce{CO}} = k_{\ce{CO}} \exp \left( \frac{-\Delta H_{\ce{CO}}}{RT} \right)
		\label{eq:K_CO_2}
	\end{equation}
	
	\colorbox{pink}{\parbox{\textwidth}{
			\textbf{\underline{IMPORTANTE}}: Para os hidrocarbonetos com $n = 1$ e $n = 2$ substitua a energia de ativação da \autoref{eq:k_2} ($E_\mathrm{HCs}$) por $E_{\mathrm{C1}}$ e $E_{\mathrm{C2}}$ (\autoref{tab:params_Mousavi_2}), respectivamente.
	}}
	
	\begin{table*}[!h]
		\centering
		\caption{Parâmetros das Equações \ref{eq:r_HCs_2}, \ref{eq:k_2} e \ref{eq:K_CO_2}.}
		\label{tab:params_Mousavi_2}
		\resizebox{1\textwidth}{!}{%
			\begin{tabular}{rrrrrr}
				\toprule
				$k_{\ce{CO}}$ (\unit{\pascal^{-1}}) & $\Delta H_{\ce{CO}}$ (\unit{\joule \per \mole}) & $A_{\mathrm{HCs}}'$ (\unit{\mole \per \second \per \gram \per \pascal^{1.75}}) & $E_{\mathrm{HCs}}$ (\unit{\joule \per \mole}) & $E_{\mathrm{C1}}$ (\unit{\joule \per \mole}) & $E_{\mathrm{C2}}$ (\unit{\joule \per \mole}) \\
				\midrule
				\num{1.398E-05} & \num{933.0} & \num{9.551E-01} & \num{147204} & \num{142602} & \num{153273} \\
				\bottomrule
			\end{tabular}
		}
	\end{table*}
	
	\subsubsection{Distribuição dos hidrocarbonetos}
	
	\begin{equation}
		y_n = (1 - \alpha) \alpha^{(n - 1)}
		\label{eq:alpha_2}
	\end{equation}
	
	$y_n$ é a fração molar do hidrocarboneto com $n$ carbonos. O $y_n$ representa exclusivamente a distribuição dos hidrocarbonetos, não estão inclusos o \ce{CO}, \ce{CO2}, \ce{H2} e \ce{H2O}.
	
	A probabilidade de crescimento de cadeia ($\alpha$) é:
	
	\begin{equation}
		\alpha = \hat{\beta_0} + \hat{\beta_1} T
	\end{equation}
	
	Onde $T$ está em Kelvin e [$\hat{\beta_0}, \hat{\beta_1}] = [\num{1.626}, \qty{-1.495e-03}{\per \kelvin}]$.
	
	\subsection{Implementação e validação no DWSIM}
	
	No DWSIM a velocidade de consumo de \ce{CO} para a reação de formação de gases foi implementada como mostra o \autoref{lst:dwsim_gases}. Para os demais componentes a implementação é mais simples, um exemplo para o querosene é dado no \autoref{lst:dwsim_querosene}. Nos \textit{scripts} abaixo $r$  é a velocidade de consumo do monóxido de carbono (é positivo mesmo, o DWSIM coloca o sinal).
	
	\colorbox{yellow}{\parbox{\textwidth}{
			\textbf{\underline{CUIDADO}}: Nos scripts abaixo os $r$ são velocidades de reação para reações isoladas, por isso não aparecem todos os termos representados na \autoref{eq:r_monox}.
	}}
	
	\begin{lstlisting}[language=Python, caption=\textit{Script} do DWSIM (Python) para implementação da velocidade de consumo de \ce{CO} na geração de gases., label={lst:dwsim_gases}]
from math import exp

R = 8.314 # J/mol K

A_HCs = 0.9550608601894232 # kmol/kg s Pa**(a+b)
E_HCs = 147204 # J/mol
E_C1 = 142602 # J/mol
E_C2 = 153273 # J/mol
k_CO = 1.397972209918762e-05 # Pa**-c
H_CO = 932.9790718460237 # J/mol
a = 0.75
b = 1.00
c = 1.00

P_H2 = R1
P_CO = R2

alfa = 1.62628872e+00 - 1.49497171e-03 * T

r_HCs_denominator = (1 + k_CO*exp(-H_CO/(R*T))*(P_CO**c))**2

# r_CO_1
y_1 = (1-alfa)*alfa**(1-1)
r_CO_1_numerator = A_HCs*exp(-E_C1/(R*T))*(P_H2**a)*(P_CO**b)
r_CO_1 = 1 * y_1 * r_CO_1_numerator/r_HCs_denominator

# r_CO_2
y_2 = (1-alfa)*alfa**(2-1)
r_CO_2_numerator = A_HCs*exp(-E_C2/(R*T))*(P_H2**a)*(P_CO**b)
r_CO_2 = 2 * y_2 * r_CO_2_numerator/r_HCs_denominator

r = r_CO_1 + r_CO_2 # kmol/kg s

# r_CO_3 e r_CO_4
r_HCs_numerator = A_HCs*exp(-E_HCs/(R*T))*(P_H2**a)*(P_CO**b)
r_HCs = r_HCs_numerator/r_HCs_denominator
for n in range(3, 5):
	y_n = (1-alfa)*alfa**(n-1)
	r += n * y_n * r_HCs # kmol/kg s
	\end{lstlisting}
	
		\begin{lstlisting}[language=Python, caption=\textit{Script} do DWSIM (Python) para implementação da velocidade de consumo de \ce{CO} na geração de querosene., label={lst:dwsim_querosene}]
from math import exp

R = 8.314 # J/mol K

A_HCs = 0.9550608601894232 # kmol/kg s Pa**(a+b)
E_HCs = 147204 # J/mol
k_CO = 1.397972209918762e-05 # Pa**-c
H_CO = 932.9790718460237 # J/mol
a = 0.75
b = 1.00
c = 1.00

P_H2 = R1
P_CO = R2

alfa = 1.62628872e+00 - 1.49497171e-03 * T


r_HCs_numerator = A_HCs*exp(-E_HCs/(R*T))*(P_H2**a)*(P_CO**b)
r_HCs_denominator = (1 + k_CO*exp(-H_CO/(R*T))*(P_CO**c))**2

r_HCs = r_HCs_numerator/r_HCs_denominator

soma_n_yn = 0

for n in range(11, 17):
	y_n = (1-alfa)*alfa**(n-1)
	soma_n_yn += n * y_n

r = soma_n_yn * r_HCs # kmol/kg s
	\end{lstlisting}
	
	\subsection{Validação do modelo de parâmetros concentrados no DWSIM}
	
	O modelo de parâmetros concentrados foi implementado no DWSIM e comparado aos dados experimentais. O resultado está representado no gráfico de paridade da \autoref{fig:parity_conc}. O preço que se paga pela simplificação é uma maior dispersão dos resíduos, o que é de se esperar, pois estamos somando as incertezas do modelo de parâmetros distribuídos (\autoref{fig:parity_dist}).
	
	\begin{figure}
		\centering
		\caption{Gráfico de paridade do modelo de parâmetros concentrados.}
		\includegraphics{parity_lumped}
		\label{fig:parity_conc}
	\end{figure}
	
	\begin{figure}
		\centering
		\caption{Gráfico de paridade do modelo de parâmetros distribuídos.}
		\includegraphics[width = \textwidth]{figure-m_out_parity_dwsim}
		\label{fig:parity_dist}
	\end{figure}
	
\end{document}